% ==================================================================
% id:       cond-026
% topic:    הסתברות מותנית ואי-תלות
% subtopic: מאורעות בלתי תלויים
% solution: full
% ==================================================================
מסקר בין המנויות והמנויים לקונצרטים של התזמורת הפילהרמונית עולים הנתונים הבאים: 60\% מהמנויים הן נשים, 40\% מהגברים המנויים גרים באזור תל אביב, ו־75\% מהמנויים הגרים באזור תל אביב הן נשים. בוחרים מנוי אחד באקראי.
\begin{enumerate}
\item מה ההסתברות שהמנוי שנבחר גר באזור תל אביב?
\hint{חשבו קודם את ההסתברות שהמנוי הוא גבר שגר באזור תל אביב, בשתי דרכים.}
\begin{solbox}
נסמן $M$ = „המנוי הוא גבר”, $T$ = „המנוי גר באזור תל אביב”. נתון $\PP(M)=0.4$, $\PP(T\mid M)=0.4$ ו־$\PP(M\mid T)=0.25$. לפי כלל הכפל $\PP(T\cap M)=0.4\cdot0.4=0.16$. מצד שני $\PP(T\cap M)=\PP(M\mid T)\PP(T)=0.25\,\PP(T)$, ולכן $\PP(T)=\frac{0.16}{0.25}=\ans{0.64}$.
\end{solbox}
\item האם המאורעות „המנוי אינו גר באזור תל אביב” ו„המנוי הוא גבר” בלתי תלויים?
\begin{solbox}
$\PP(T^C\cap M)=\PP(M)-\PP(T\cap M)=0.4-0.16=0.24$, ואילו $\PP(T^C)\PP(M)=0.36\cdot0.4=0.144$. לכן המאורעות \(\ans{\text{תלויים}}\).

אפשר גם כך: $T^C$ ו־$M$ בלתי תלויים אם ורק אם $T$ ו־$M$ בלתי תלויים, ו־$\PP(T\mid M)=0.4\ne0.64=\PP(T)$.
\end{solbox}
\end{enumerate}
